% ==============================================================================
% CAPÍTULO 1 - INTRODUÇÃO
% ==============================================================================

\chapter{Introdução}
\label{cap:introducao}

A transformação digital tem impactado de forma crescente os processos internos das empresas industriais,
tornando a automação de documentos e a integração entre sistemas operacionais cada vez mais estratégicas
para a competitividade. Nesse contexto, o desenvolvimento de software empresarial, em especial sistemas
de gestão customizados, emerge como uma disciplina central da Engenharia de Software, exigindo dos
profissionais habilidades que vão desde a modelagem de dados até a geração automatizada de documentos com
alto padrão visual.

A BIOTECNO Indústria e Comércio Ltda., empresa fabricante de equipamentos de refrigeração científica
localizada em Santa Rosa (RS), opera com um sistema web interno denominado Bioweb, desenvolvido em
Django. Dentro do módulo comercial do Bioweb, a geração de propostas técnicas e comerciais é uma das
atividades de maior impacto direto no relacionamento com clientes, pois o documento produzido representa
a empresa perante o mercado.

Até o início deste estágio, os orçamentos em PDF eram gerados por meio da biblioteca PyLaTeX, que
construía o documento de forma programática e estática, ou seja, qualquer alteração no layout da
proposta exigia modificação direta no código-fonte Python. Esse modelo limitava a capacidade de
atualização visual das propostas e criava dependência técnica para ajustes que, em essência, são decisões
de negócio e identidade visual. Paralelamente, a empresa contratou uma empresa de design para a criação
de um novo layout de proposta comercial, o qual precisava ser integrado ao sistema existente de forma
flexível e sustentável.

Este relatório descreve as atividades desenvolvidas ao longo do estágio curricular obrigatório realizado
na BIOTECNO, com foco na concepção e implementação de um sistema de geração automatizada de propostas
comerciais em PDF. O trabalho envolveu a codificação do template LaTeX a partir do layout entregue pela
empresa de design, a modelagem de um novo conjunto de entidades no banco de dados, e o desenvolvimento
da lógica de geração de documentos na camada de \textit{views} do Django com ênfase em uma arquitetura
que permita a troca de layouts diretamente pelo painel administrativo, sem necessidade de alterações no
código.

%-------------------------------------------------
% Justificativa
%-------------------------------------------------
\section{Justificativa}
\label{sec:justificativa}

A prática de estágio supervisionado constitui etapa fundamental na formação do profissional de Ciência da
Computação, pois permite a aplicação concreta dos conhecimentos construídos ao longo do curso em
problemas reais de complexidade e escala incompatíveis com o ambiente acadêmico isolado. A exposição a
decisões de arquitetura de software com impacto direto em processos de negócio desenvolve no estagiário
uma visão sistêmica que complementa a formação teórica.

Do ponto de vista acadêmico, o projeto demandou a aplicação integrada de conteúdos de Engenharia de
Software (modelagem de requisitos, padrões de projeto), Banco de Dados (modelagem relacional, migrações),
Desenvolvimento Web (arquitetura MVC/MTV, \textit{views} baseadas em função) e Sistemas de Composição
Tipográfica (\LaTeX{}), demonstrando a interdisciplinaridade inerente ao desenvolvimento de sistemas
reais. A necessidade de integrar dois sistemas com gramáticas conflitantes --- o motor de templates do
Django e a sintaxe do \LaTeX{} --- e de tomar decisões arquiteturais com impacto direto na
manutenibilidade do sistema são exemplos de desafios que a formação acadêmica prepara teoricamente,
mas que só se consolidam no contato com um ambiente produtivo real.

%-------------------------------------------------
% Objetivos
%-------------------------------------------------
\section{Objetivos}
\label{sec:objetivos}

\subsection{Objetivo Geral}
\label{subsec:objective-geral}

Desenvolver um módulo de geração automatizada de propostas técnicas e comerciais em PDF, integrado ao
sistema Django da BIOTECNO, com suporte a layouts intercambiáveis gerenciáveis pelo painel administrativo,
substituindo a abordagem estática anterior baseada na biblioteca PyLaTeX.

\subsection{Objetivos Específicos}
\label{subsec:objetivos-especificos}

\begin{alineas}
    \item Analisar o processo vigente de geração de orçamentos e levantar os requisitos funcionais e
    não funcionais do novo sistema;
    \item Traduzir o layout visual elaborado pela empresa de design para um template \LaTeX{} funcional,
    parametrizável via sistema de templates do Django;
    \item Modelar e implementar o modelo \texttt{LayoutOrcamento} no banco de dados, incluindo suporte
    ao upload de templates compactados e ao mecanismo de layout ativo único;
    \item Desenvolver as funções de resolução de diretório de layout, montagem de contexto e geração de
    PDF na camada de \textit{views} do Django, integrando o pipeline completo do banco de dados à
    entrega do documento ao navegador;
    \item Documentar o processo de desenvolvimento e os procedimentos necessários para manutenção
    futura da solução.
\end{alineas}

%-------------------------------------------------
% Estrutura do Trabalho
%-------------------------------------------------
\section{Estrutura do Trabalho}
\label{sec:estrutura}

Para melhor compreensão, este relatório está organizado em capítulos sequenciais. Além desta Introdução,
o \autoref{ch:apresentacao-empresa} apresenta o perfil da BIOTECNO Indústria e Comércio Ltda. e o
ambiente tecnológico em que o estágio foi realizado. O \autoref{cap:metodologia} descreve a sistemática
de trabalho adotada, as ferramentas utilizadas e a forma como as demandas foram gerenciadas ao longo do
período. O \autoref{ch:revisao-bibliografica} fundamenta tecnicamente o projeto, discutindo os principais
conceitos e tecnologias empregados. O \autoref{cap:desenvolvimento} detalha o processo completo de
análise, modelagem e implementação do sistema de geração de orçamentos. Por fim, o
\autoref{cap:consideracoes} apresenta as conclusões e perspectivas de trabalhos futuros.